%Este es documento para el libro del proyecto
%Entrega: Anteproyecto - Parte 3
%Compilador en Overleaf: Menu > Compiler > LaTeX (latex + dvips + ps2pdf)
 \documentclass[11pt,openany,oneside,letterpaper]{book}
 \usepackage[spanish]{babel}
 %\usepackage[latin1]{inputenc}
 \usepackage[utf8]{inputenc}

%%Dimensiones de página
 \usepackage[hmargin={3.0cm,2.5cm},vmargin={3.0cm,2.5cm}]{geometry}
 \usepackage[spanish]{layout}
 \marginparwidth=2.0cm
 %\layout %Para mirar como esta configurada la página

%Sangria y espacio entre párrafos
 %\parindent=0pt
 \parskip=10pt

%Paquetes empleados
 \usepackage{array,amsmath,amssymb,amsfonts,pifont,verbatim}
 \usepackage{dsfont}

%Paquetes para el formato del libro
 \usepackage{fancyhdr}
 \usepackage[]{fncychap}

%Paquetes para las gráficas
 \usepackage{pstricks,pst-all}

%Paquete de manejo gráfico
 \usepackage[dvips]{graphicx}
 %Carpeta de imágenes (en Overleaf: crear la carpeta multimedia_latex y subir allí EscudoUD.eps)
 \graphicspath{{multimedia_latex/}}

%Paquete para insertar mejor los objetos flotantes
 \usepackage{float}

%Paquetes agregados para tablas, figuras vectoriales y enlaces
 \usepackage{booktabs,tabularx,longtable,multirow}
 \usepackage{url}
 \usepackage{tikz}
 \usetikzlibrary{babel,positioning,arrows.meta,shapes.geometric,fit,calc,patterns}
 \usepackage{pgfplots}
 \pgfplotsset{compat=1.16}

%Sección con título corto solo en el encabezado de página (el índice conserva el título completo)
 \newcommand{\seccioncorta}[2]{{\renewcommand{\sectionmark}[1]{}\section{#2}}\markright{\MakeUppercase{\thesection.\ #1}}}
%Redefinition de títulos y rótulos
 \renewcommand{\listtablename}{\'Indice de tablas}
 \renewcommand{\tablename}{Tabla}
 \addto\captionsspanish{\renewcommand{\listtablename}{\'Indice de tablas}\renewcommand{\tablename}{Tabla}}

%Bibliografía en formato APA: sin etiquetas numéricas y con sangría francesa
 \makeatletter
 \renewenvironment{thebibliography}[1]
  {\chapter*{\bibname}\addcontentsline{toc}{chapter}{\bibname}%
   \def\@biblabel##1{}\list{}{\usecounter{enumiv}\setlength{\leftmargin}{1.27cm}\setlength{\itemindent}{-1.27cm}%
   \setlength{\itemsep}{4pt}\setlength{\parsep}{0pt}}\sloppy}
  {\endlist}
 \makeatother

\begin{document}
%==================== Inicio de PRA0Portada.tex ====================
%PÁGINA DE LA PORTADA
\thispagestyle{empty}
\begin{center}
\includegraphics[scale=0.25]{EscudoUD.eps}\\
\vspace*{1.5cm}

{\large \bf Evaluación del impacto de la calidad de datos en la precisión de resultados analíticos: un enfoque centrado en duplicados exactos}\\
\vspace*{1.2cm}

{\Large \bf Anteproyecto - Parte 3}\\
\vspace*{2.5cm}

{\large Integrantes:}\\[0.3cm]
David Martínez Rincón - 20222020167\\
Juan David Pulido Rodríguez - 20202020063\\
Santiago Alfonso Arias - 20222020093\\
\vspace*{2cm}

Docente: Helbert Eduardo Espitia Cuchango\\
\vspace*{2.5cm}

{\large METODOLOGÍA DE LA INVESTIGACIÓN\\
UNIVERSIDAD DISTRITAL FRANCISCO JOSÉ DE CALDAS\\
FACULTAD DE INGENIERÍA\\
INGENIERÍA DE SISTEMAS\\
BOGOTÁ D.C., COLOMBIA\\
2026}
\end{center}
\newpage
%==================== Fin de PRA0Portada.tex ====================

%Numeración inicial
 \frontmatter
 \tableofcontents
 %\listoffigures
 %\listoftables

%Numeración principal
 \mainmatter
 \setcounter{chapter}{3}
%==================== Inicio de PRB7Tiempo.tex ====================
\chapter{Tiempo planeado}

El proyecto se desarrolla a lo largo de las 17 semanas del semestre, del 3 de agosto al 27 de noviembre de 2026. Las primeras diez ya se cumplieron: en ellas se definió el tema, se hizo la revisión documental, se presentaron las entregas del curso en las fechas de cierre que aparecen en la Tabla~\ref{tab:cronograma} y se escogieron los estudios ancla. En la semana 10 reevaluamos el ancla inicial y, después de verificar de forma independiente varios candidatos, en la semana 11 la cambiamos por Many Labs~1 (Klein et al., 2014), con el conjunto de datos Saber~11/Saber~Pro (Delahoz-Dominguez et al., 2020) como segundo caso. Las seis semanas restantes se dedican a ejecutar la prueba de concepto, al experimento completo, a la redacción del documento final y a la presentación del proyecto a finales de noviembre. Las fechas descuentan los festivos del 7 y el 17 de agosto, el 12 de octubre y el 2 y el 16 de noviembre. La Figura~\ref{fig:gantt} resume la distribución de las etapas en el tiempo.

\begin{figure}[H]
\centering
\begin{tikzpicture}[font=\scriptsize]
\node at (0.280,0.45) {1};
\node at (0.840,0.45) {2};
\node at (1.400,0.45) {3};
\node at (1.960,0.45) {4};
\node at (2.520,0.45) {5};
\node at (3.080,0.45) {6};
\node at (3.640,0.45) {7};
\node at (4.200,0.45) {8};
\node at (4.760,0.45) {9};
\node at (5.320,0.45) {10};
\node at (5.880,0.45) {11};
\node at (6.440,0.45) {12};
\node at (7.000,0.45) {13};
\node at (7.560,0.45) {14};
\node at (8.120,0.45) {15};
\node at (8.680,0.45) {16};
\node at (9.240,0.45) {17};
\draw[gray!30] (0.000,0.2) -- (0.000,-5.48);
\draw[gray!30] (0.560,0.2) -- (0.560,-5.48);
\draw[gray!30] (1.120,0.2) -- (1.120,-5.48);
\draw[gray!30] (1.680,0.2) -- (1.680,-5.48);
\draw[gray!30] (2.240,0.2) -- (2.240,-5.48);
\draw[gray!30] (2.800,0.2) -- (2.800,-5.48);
\draw[gray!30] (3.360,0.2) -- (3.360,-5.48);
\draw[gray!30] (3.920,0.2) -- (3.920,-5.48);
\draw[gray!30] (4.480,0.2) -- (4.480,-5.48);
\draw[gray!30] (5.040,0.2) -- (5.040,-5.48);
\draw[gray!30] (5.600,0.2) -- (5.600,-5.48);
\draw[gray!30] (6.160,0.2) -- (6.160,-5.48);
\draw[gray!30] (6.720,0.2) -- (6.720,-5.48);
\draw[gray!30] (7.280,0.2) -- (7.280,-5.48);
\draw[gray!30] (7.840,0.2) -- (7.840,-5.48);
\draw[gray!30] (8.400,0.2) -- (8.400,-5.48);
\draw[gray!30] (8.960,0.2) -- (8.960,-5.48);
\draw[gray!30] (9.520,0.2) -- (9.520,-5.48);
\node[anchor=east, font=\scriptsize\bfseries] at (-0.15,0.45) {Semana};
\node[anchor=east, align=right] at (-0.15,-0.310) {1. Problema y revisión documental};
\draw[fill=gray!65] (0.030,-0.510) rectangle (3.330,-0.110);
\node[anchor=east, align=right] at (-0.15,-0.930) {2. Selección del ancla y anteproyecto};
\draw[fill=gray!65] (3.390,-1.130) rectangle (5.570,-0.730);
\draw[pattern=north east lines, pattern color=gray!80] (5.630,-1.130) rectangle (6.130,-0.730);
\node[anchor=east, align=right] at (-0.15,-1.550) {3. Prueba de concepto (Saber)};
\draw[fill=white] (6.190,-1.750) rectangle (7.250,-1.350);
\node[anchor=east, align=right] at (-0.15,-2.170) {4. Preparación de Many Labs~1};
\draw[fill=white] (7.310,-2.370) rectangle (7.810,-1.970);
\node[anchor=east, align=right] at (-0.15,-2.790) {5. Análisis inferencial};
\draw[fill=white] (7.310,-2.990) rectangle (7.810,-2.590);
\node[anchor=east, align=right] at (-0.15,-3.410) {6. Tamaños de efecto y consolidación};
\draw[fill=white] (7.870,-3.610) rectangle (8.370,-3.210);
\node[anchor=east, align=right] at (-0.15,-4.030) {7. Comparación y análisis};
\draw[fill=white] (8.430,-4.230) rectangle (8.930,-3.830);
\node[anchor=east, align=right] at (-0.15,-4.650) {8. Documento final y presentación};
\draw[fill=white] (8.430,-4.850) rectangle (9.490,-4.450);
\node[anchor=east, font=\scriptsize\bfseries] at (-0.15,-5.270) {Entregas del curso};
\node[diamond, draw, fill=black, inner sep=1.6pt] at (1.960,-5.270) {};
\node[diamond, draw, fill=black, inner sep=1.6pt] at (2.520,-5.270) {};
\node[diamond, draw, fill=black, inner sep=1.6pt] at (3.080,-5.270) {};
\node[diamond, draw, fill=black, inner sep=1.6pt] at (4.200,-5.270) {};
\node[diamond, draw, fill=black, inner sep=1.6pt] at (4.760,-5.270) {};
\node[diamond, draw, fill=black, inner sep=1.6pt] at (9.240,-5.270) {};
\draw[red!70!black, thick, dashed] (5.824,0.2) -- (5.824,-5.48) node[below, text=red!70!black] {14 de octubre};
\draw[fill=gray!65] (0,-6.33) rectangle ++(0.45,0.28); \node[anchor=west] at (0.5,-6.19) {realizado};
\draw[pattern=north east lines, pattern color=gray!80] (2.2,-6.33) rectangle ++(0.45,0.28); \node[anchor=west] at (2.7,-6.19) {en curso};
\draw (4.3,-6.33) rectangle ++(0.45,0.28); \node[anchor=west] at (4.8,-6.19) {pendiente};
\node[diamond, draw, fill=black, inner sep=1.6pt] at (6.6,-6.19) {}; \node[anchor=west] at (6.75,-6.19) {entrega};
\end{tikzpicture}
\caption{Diagrama de Gantt del proyecto por etapas, con las entregas del curso. La línea punteada marca la fecha de esta entrega.}
\label{fig:gantt}
\end{figure}

{\small
\setlength{\tabcolsep}{4pt}
\begin{longtable}{c >{\raggedright\arraybackslash}p{1.8cm} >{\raggedright\arraybackslash}p{4.6cm} >{\raggedright\arraybackslash}p{4.3cm} >{\raggedright\arraybackslash}p{1.5cm}}
\caption{Cronograma del proyecto}\label{tab:cronograma}\\
\toprule
\textbf{Sem.} & \textbf{Fechas} & \textbf{Actividad principal} & \textbf{Resultado o entrega} & \textbf{Estado} \\
\midrule
\endfirsthead
\toprule
\textbf{Sem.} & \textbf{Fechas} & \textbf{Actividad principal} & \textbf{Resultado o entrega} & \textbf{Estado} \\
\midrule
\endhead
\bottomrule
\endfoot
1--3 & 3 -- 21 ago & Definición del tema, problema y pregunta de investigación; búsqueda de antecedentes en revistas indexadas & Propuesta de tema y primeros antecedentes & Realizado \\
4 & 24 -- 28 ago & Matriz de revisión documental y consolidación del tema & Entregas: revisión documental (25 ago) y tema del proyecto (28 ago) & Realizado \\
5 & 31 ago -- 4 sep & Organización de los antecedentes por ejes temáticos & Entrega: mapa mental (3 sep) & Realizado \\
6 & 7 -- 11 sep & Redacción del estado del arte & Entrega: estado del arte (8 sep) & Realizado \\
7 & 14 -- 18 sep & Búsqueda y evaluación de candidatos a estudio ancla; diseño de la prueba de concepto & Candidatos evaluados con seis criterios & Realizado \\
8 & 21 -- 25 sep & Auditoría del código del estudio ancla y decisión inicial & Entrega: prueba de concepto (21 sep); ancla inicial: Patel (2025) & Realizado \\
9 & 28 sep -- 2 oct & Redacción del anteproyecto & Entregas: anteproyecto partes 1 y 2 (30 sep) y parte 3 (1 oct) & Realizado \\
10 & 5 -- 9 oct & Reevaluación del ancla: búsqueda de candidatos revisados por pares, descarga de sus datos y réplica independiente de sus resultados & Informe de candidatos; Many Labs~1 y Saber~11/Saber~Pro verificados & Realizado \\
11 & 13 -- 16 oct & Cambio del ancla a Many Labs~1, con Saber~11/Saber~Pro como segundo caso; actualización de las entregas en el formato \LaTeX{} del curso & Estado del arte, prueba de concepto y anteproyecto actualizados & En curso \\
12 & 19 -- 23 oct & Prueba de concepto: carga de Saber~11/Saber~Pro en PostgreSQL, línea base, scripts de inyección y generación de las 160 versiones & Línea base verificada y versiones con duplicados & Pendiente \\
13 & 26 -- 30 oct & Prueba de concepto: consultas sobre todas las versiones, deduplicación, métricas, gráficas e informe & Informe de la prueba de concepto y ajustes al diseño & Pendiente \\
14 & 3 -- 6 nov & Many Labs~1: preparación de los datos, réplica de los cinco efectos en Python, generación de versiones y pruebas chi-cuadrado y $t$ & Línea base replicada y estadísticos de todas las versiones & Pendiente \\
15 & 9 -- 13 nov & Many Labs~1: tamaños de efecto (V de Cramér y $d$ de Cohen); consolidación de los resultados de los dos casos & Tabla consolidada de resultados descriptivos e inferenciales & Pendiente \\
16 & 17 -- 20 nov & Comparación, análisis de resultados y redacción del documento final en el formato de tesis del curso & Documento final en borrador & Pendiente \\
17 & 23 -- 27 nov & Revisión final, preparación de la presentación y presentación del proyecto & Documento final y presentación & Pendiente \\
\end{longtable}
}
%==================== Fin de PRB7Tiempo.tex ====================
%==================== Inicio de PRB8Limitaciones.tex ====================
\chapter{Limitaciones}

Desde el inicio identificamos varias limitaciones. Las dejamos por escrito junto con la forma en que pensamos manejarlas, para que los resultados se interpreten en su justa medida.

\textbf{Los estudios ancla son pequeños y de dominios específicos.} Many Labs~1 tiene 6\,344 participantes y sus efectos vienen de la psicología social; Saber~11/Saber~Pro tiene 12\,411 estudiantes de ingeniería. Son conjuntos revisados por pares, públicos y reproducibles, pero no representan las bases de datos de millones de registros que maneja una organización grande. Lo que encontremos describe el comportamiento de estos datos, estos indicadores y estas pruebas, y no debe extenderse a cualquier base de datos sin nuevas pruebas. Como el efecto de los duplicados depende de su proporción y no del tamaño absoluto de la base, esperamos que los patrones se mantengan, pero eso es algo que este proyecto no comprueba.

\textbf{Los datos de Many Labs~1 exigen una preparación cuidadosa.} El archivo tiene 382 columnas, varias respuestas vacías vienen como espacios en blanco y no como valores nulos, y el análisis original se publicó en SPSS. Por eso trataremos los espacios en blanco como faltantes, reimplementaremos el análisis en Python y solo trabajaremos con los cinco efectos cuyos resultados publicados reproducimos con exactitud.

\textbf{Saber~11/Saber~Pro solo tiene resultados descriptivos.} El artículo de datos reporta promedios, desviaciones y porcentajes, pero ninguna prueba estadística. Por eso cada caso cubre un nivel de análisis distinto: Saber el descriptivo y Many Labs el inferencial. El archivo, además, viene en Excel y trae una columna vacía que se descarta al cargarlo.

\textbf{Solo se estudian duplicados exactos.} Quedan por fuera los registros parecidos pero no idénticos y los demás problemas de calidad: datos faltantes, valores incorrectos, inconsistencias entre campos y errores de formato.

\textbf{Los duplicados son sintéticos.} Los introducimos nosotros, en porcentajes de 1\,\%, 5\,\%, 10\,\% y 20\,\% y con tres tipos de esquema (uniforme, concentrado en un grupo natural y concentrado según el resultado). Esto permite controlar exactamente cuánto y dónde se duplica, pero no cubre todos los patrones de duplicación que pueden aparecer en un sistema real. Los esquemas concentrados buscan acercarse a fallas típicas, como una carga que se ejecuta dos veces para una sola fuente. En algunos esquemas el grupo es más pequeño que el número de filas que hay que agregar con 20\,\%, así que esos casos se reportarán aparte.

\textbf{Los valores $p$ de la línea base son prácticamente nulos.} Los cinco efectos de Many Labs~1 que replicamos son muy significativos, así que es poco probable que los duplicados los hagan cruzar el umbral de 0,05 en cualquiera de las dos direcciones. Por eso no mediremos si cambia la decisión de la prueba, sino cuánto cambian la magnitud del estadístico y el tamaño de efecto, que es lo que determina qué tan fuerte parece un hallazgo.

\textbf{El tiempo es corto.} Quedan seis semanas, y el cambio de ancla en la semana 11 consumió parte del margen que teníamos. Lo compensa que los dos conjuntos de datos nuevos son pequeños y que ya verificamos sus líneas base, así que la carga y la réplica son más rápidas que con el ancla anterior. Si el tiempo no alcanza, priorizaremos el efecto de encuadre de Many Labs~1 y los indicadores de Saber que ya probamos en la exploración preliminar.

\textbf{Reproducibilidad y dependencia de las herramientas.} La eliminación de duplicados exactos con \texttt{SELECT DISTINCT} funciona igual en cualquier gestor relacional, así que los resultados no dependen de PostgreSQL. Lo que sí depende de la herramienta es la generación de números aleatorios para escoger qué registros se duplican, que en PostgreSQL se controla con \texttt{setseed}. Para garantizar que el experimento se pueda repetir, registraremos todas las semillas y guardaremos las versiones generadas.
%==================== Fin de PRB8Limitaciones.tex ====================
%==================== Inicio de PRB9Plan.tex ====================
\chapter{Plan de trabajo}

El trabajo está organizado en ocho etapas que van desde la definición del problema hasta la presentación final (Figura~\ref{fig:etapas}). Las dos primeras ya están completas o a punto de cerrarse. La prueba de concepto cumple un doble papel: ensaya a escala el ciclo completo con Saber~11/Saber~Pro y, como ese conjunto de datos se usa entero, sus resultados son directamente los del caso descriptivo. Los scripts que se validen ahí se reutilizan después con Many Labs~1.

\begin{figure}[H]
\centering
\begin{tikzpicture}[
  et/.style={draw, rounded corners=3pt, align=center, text width=3.05cm, minimum height=2.0cm, font=\scriptsize, inner sep=3pt, execute at begin node={\hyphenpenalty=10000\relax}},
  hecho/.style={et, fill=gray!25},
  >={Stealth[length=2.3mm]}]
\node[hecho] (e1) at (0,0) {\textbf{1. Problema y revisión}\\ sem. 1--6\\[2pt] \emph{estado del arte}};
\node[hecho] (e2) at (3.75,0) {\textbf{2. Ancla y anteproyecto}\\ sem. 7--11\\[2pt] \emph{Many Labs~1 y Saber verificados}};
\node[et] (e3) at (7.5,0) {\textbf{3. Prueba de concepto}\\ sem. 12--13\\[2pt] \emph{scripts validados con Saber}};
\node[et] (e4) at (11.25,0) {\textbf{4. Preparación de Many Labs~1}\\ sem. 14\\[2pt] \emph{línea base replicada}};
\node[et] (e5) at (11.25,-2.7) {\textbf{5. Análisis inferencial}\\ sem. 14\\[2pt] \emph{$\chi^2$, $t$ y valores $p$}};
\node[et] (e6) at (7.5,-2.7) {\textbf{6. Tamaños de efecto y consolidación}\\ sem. 15\\[2pt] \emph{tabla consolidada}};
\node[et] (e7) at (3.75,-2.7) {\textbf{7. Comparación}\\ sem. 16\\[2pt] \emph{tablas y gráficas}};
\node[et] (e8) at (0,-2.7) {\textbf{8. Documento y presentación}\\ sem. 16--17\\[2pt] \emph{documento final}};
\draw[->] (e1) -- (e2); \draw[->] (e2) -- (e3); \draw[->] (e3) -- (e4);
\draw[->] (e4) -- (e5); \draw[->] (e5) -- (e6); \draw[->] (e6) -- (e7); \draw[->] (e7) -- (e8);
\draw[->, dashed] (e3.south west) ++(0.3,0) -- node[left, font=\scriptsize, align=right] {resultados\\ descriptivos} ([xshift=0.3cm]e6.north west);
\end{tikzpicture}
\caption{Etapas del plan de trabajo, con sus semanas y el producto principal de cada una. Las etapas sombreadas ya se completaron o están por cerrarse; la línea punteada indica que los resultados de la prueba de concepto pasan directamente a la consolidación.}
\label{fig:etapas}
\end{figure}

\seccioncorta{Etapa 1}{Etapa 1. Definición del problema y revisión documental (semanas 1 a 6, realizada)}

Se definió el problema, la pregunta y los objetivos, y se hizo una revisión documental de veinte antecedentes en revistas colombianas y latinoamericanas indexadas, complementada con referentes internacionales. La revisión mostró que el efecto de los duplicados solo se ha medido sobre la exactitud de clasificadores y que falta evidencia sobre indicadores descriptivos, pruebas estadísticas y conclusiones de estudios reales.

\textbf{Producto:} tema del proyecto, matriz de revisión documental, mapa mental y estado del arte, entregados entre el 25 de agosto y el 8 de septiembre.

\seccioncorta{Etapa 2}{Etapa 2. Selección de los estudios ancla y anteproyecto (semanas 7 a 11)}

En las semanas 7 y 8 se evaluaron cinco candidatos con seis criterios: datos públicos, metodología documentada, análisis replicable en el tiempo disponible, reputación de la fuente, tamaño manejable y un significado claro de ``duplicado exacto''. Se escogió inicialmente el estudio de Patel (2025), y con él se presentaron la prueba de concepto y el anteproyecto. En la semana 10 reevaluamos esa decisión, porque se trataba de un preprint sin revisión por pares y su código no coincidía del todo con el artículo. Hicimos una nueva búsqueda centrada en estudios revisados por pares, descargamos los datos de los candidatos más fuertes y replicamos sus resultados por nuestra cuenta. Many Labs~1 (Klein et al., 2014) y Saber~11/Saber~Pro (Delahoz-Dominguez et al., 2020) fueron los que mejor cumplieron los criterios, y sus resultados publicados los reprodujimos con exactitud, así que en la semana 11 los adoptamos como anclas y actualizamos las entregas.

\textbf{Producto:} estudios ancla definidos y verificados, propuesta de la prueba de concepto y anteproyecto, actualizados en el formato \LaTeX{} del curso.

\seccioncorta{Etapa 3}{Etapa 3. Ejecución de la prueba de concepto (semanas 12 y 13)}

Sobre los 12\,411 registros de Saber~11/Saber~Pro se ejecuta el ciclo completo: línea base, inyección de duplicados en cuatro niveles y cuatro esquemas con 10 semillas por combinación, ocho indicadores calculados con consultas SQL idénticas, deduplicación y comparación.

\textbf{Producto:} informe de la prueba de concepto, resultados del caso descriptivo y scripts validados.

\seccioncorta{Etapa 4}{Etapa 4. Preparación de los datos de Many Labs~1 (semana 14)}

Se lee el archivo depurado de Many Labs~1, se tratan los espacios en blanco como valores faltantes, se carga en PostgreSQL y se verifica que no haya duplicados exactos. Luego se reimplementan en Python las pruebas de los cinco efectos seleccionados y se comprueba que reproducen los valores publicados. Con los mismos scripts de la prueba de concepto se generan las versiones con 1\,\%, 5\,\%, 10\,\% y 20\,\% de duplicados en los esquemas uniforme, concentrado en el sitio de recolección más grande y concentrado según el resultado.

\textbf{Producto:} línea base replicada y versiones con duplicados verificadas.

\seccioncorta{Etapa 5}{Etapa 5. Análisis inferencial (semana 14)}

Sobre cada versión y sobre su versión deduplicada se ejecutan las tres pruebas chi-cuadrado y las dos pruebas $t$, y se registran el estadístico, los grados de libertad y el valor $p$.

\textbf{Producto:} tabla de resultados inferenciales de todas las versiones.

\seccioncorta{Etapa 6}{Etapa 6. Tamaños de efecto y consolidación (semana 15)}

Se calculan la V de Cramér para las pruebas chi-cuadrado y la $d$ de Cohen para las pruebas $t$ en todas las versiones, y se reúnen en una sola tabla los resultados de los dos casos: los descriptivos de Saber y los inferenciales de Many Labs.

\textbf{Producto:} tabla consolidada de resultados.

\seccioncorta{Etapa 7}{Etapa 7. Comparación y análisis de resultados (semana 16)}

Se calculan el error relativo, las diferencias en puntos porcentuales, la tau de Kendall para los rankings, la razón de cada estadístico frente a la línea base y el cambio en los tamaños de efecto. Se grafica cada resultado frente al nivel de duplicación, separando los esquemas, para ver qué resultados son más sensibles y a partir de qué nivel la distorsión se vuelve relevante. Los hallazgos se contrastan con las hipótesis y con la literatura revisada.

\textbf{Producto:} análisis de resultados con tablas y gráficas.

\seccioncorta{Etapa 8}{Etapa 8. Documento final y presentación (semanas 16 y 17)}

Se redacta el documento final en el formato de tesis del curso, con la metodología, los resultados, las conclusiones, las limitaciones encontradas y el trabajo futuro, y se prepara la presentación del proyecto para la semana del 23 al 27 de noviembre.

\textbf{Producto:} documento final y presentación del proyecto.
%==================== Fin de PRB9Plan.tex ====================

%Numeración final
 \backmatter

\begin{thebibliography}{xx}
%==================== Inicio de Bibliografia.tex ====================
%Bibliografía en formato APA 7

\bibitem{delahoz2020} Delahoz-Dominguez, E., Zuluaga, R., y Fontalvo-Herrera, T. (2020). Dataset of academic performance evolution for engineering students. \emph{Data in Brief, 30}, 105537. \url{https://doi.org/10.1016/j.dib.2020.105537}

\bibitem{klein2014} Klein, R. A., Ratliff, K. A., Vianello, M., Adams, R. B., Jr., Bahník, Š., Bernstein, M. J., Bocian, K., Brandt, M. J., Brooks, B., Brumbaugh, C. C., Cemalcilar, Z., Chandler, J., Cheong, W., Davis, W. E., Devos, T., Eisner, M., Frankowska, N., Furrow, D., Galliani, E. M., \ldots{} Nosek, B. A. (2014). Investigating variation in replicability: A ``many labs'' replication project. \emph{Social Psychology, 45}(3), 142--152. \url{https://doi.org/10.1027/1864-9335/a000178}

\bibitem{patel2025} Patel, U. (2025). \emph{Identifying determinants of cesarean births: A machine learning framework for maternal health research} [Preprint]. Research Square. \url{https://doi.org/10.21203/rs.3.rs-7320875/v1}
%==================== Fin de Bibliografia.tex ====================
\end{thebibliography}
\end{document}
